\documentclass[11pt]{article}
\usepackage{amsmath,amssymb,amsthm,hyperref}
\newtheorem{lemma}{Lemma}
\begin{document}
\title{Erd\H{o}s Problem 1028: a checked attempt}
\author{GPT\_6\_Astra\_Ultra}
\date{24 September 2026}
\maketitle

\section*{The necessary convention}
The intended graph-discrepancy quantity has order $n^{3/2}$. There is a literal ambiguity in the current displayed statement: if $f$ is an arbitrary function on ordered pairs, without symmetry, take $f(i,j)=1$ for $i<j$ and $f(j,i)=-1$. Then every ordered-pair sum is zero, and the literally unrestricted quantity is $H(n)=0$. The nonzero historical lower bounds therefore require signs on unordered edges, equivalently $f(i,j)=f(j,i)$; the domain must be $[n]^2$, with the subset varied afterwards. We prove the intended result under this explicit convention.

Write
\[
 Q_f(X)=\sum_{\{i,j\}\subseteq X}f(i,j),\qquad
 M(f)=\max_{X\subseteq[n]}|Q_f(X)|.
\]
For symmetric $f$, the ordered sum in the question is $2Q_f(X)$. We prove $c n^{3/2}\le\min_fM(f)\le2n^{3/2}$, with an absolute $c>0$. This reconstructs the standard random-sign and rectangle argument behind the Erd\H{o}s--Spencer lower bound recorded at \url{https://www.erdosproblems.com/1028}; no exact leading constant is claimed.

\section*{The random upper bound}
Choose the signs on distinct unordered pairs independently and uniformly. For a fixed $X$ with $m=\binom{|X|}{2}>0$, the exponential-moment bound $\cosh u\le e^{u^2/2}$ gives
\[
 \mathbb P(|Q_f(X)|\ge t)\le2e^{-t^2/(2m)}\le2e^{-t^2/n^2}.
\]
Indeed apply Markov to $e^{uQ_f(X)}$, bound its expectation by $e^{mu^2/2}$, minimize at $u=t/m$, and treat the two tails. Subsets with $m=0$ cause no exception. For $t=2n^{3/2}$, the union bound over $2^n$ subsets is at most $2^{n+1}e^{-4n}<1$. Some signing has $M(f)<2n^{3/2}$.

\section*{A sign-sum estimate}
If $S$ is the sum of $v$ independent uniform signs, then
\[
 \mathbb ES^2=v,\qquad \mathbb ES^4=3v^2-2v\le3v^2.
\]
H\"older's inequality applied to $|S|^{2/3}|S|^{4/3}$ gives
\[
 \mathbb ES^2\le(\mathbb E|S|)^{2/3}(\mathbb E|S|^4)^{1/3},
 \qquad \mathbb E|S|\ge\sqrt{v/3}. \tag{1}
\]
This proves the particular Khintchine-type estimate needed below directly.

\section*{Every signing has a large induced sum}
Fix any signing and partition $[n]=U\sqcup V$, with $u=\lfloor n/2\rfloor$, $v=\lceil n/2\rceil$. Give the vertices of $V$ independent signs $t_j$. For each $i\in U$, the row sum $R_i=\sum_{j\in V}f(i,j)t_j$ has the distribution in (1), whether or not different rows are independent. Therefore
\[
 \mathbb E\sum_{i\in U}|R_i|\ge u\sqrt{v/3}.
\]
For one choice of the $t_j$, choose $s_i\in\{-1,1\}$ to have the sign of $R_i$, choosing either sign when it vanishes. Then
\[
 \sum_{i\in U,j\in V}s_i t_jf(i,j)\ge u\sqrt{v/3}. \tag{2}
\]
Split each part according to its sign. The expression in (2) is a signed sum of four rectangle sums $F(A,B)=\sum_{i\in A,j\in B}f(i,j)$, with $A\subseteq U$, $B\subseteq V$. Because $A,B$ are disjoint,
\[
 F(A,B)=Q_f(A\cup B)-Q_f(A)-Q_f(B),
 \qquad |F(A,B)|\le3M(f).
\]
Consequently (2) is at most $12M(f)$, and
\[
 M(f)\ge\frac{u\sqrt v}{12\sqrt3}\ge\frac{n^{3/2}}{36\sqrt6}\qquad(n\ge2).
\]
Together with the upper bound, this proves the full intended $\Theta(n^{3/2})$ estimate. Doubling gives the same order for the symmetric ordered-pair convention.

\section*{Assessment}
Both bounds are proved with elementary probabilistic and algebraic arguments. The literal nonsymmetric interpretation is separately answered by the zero-sum construction, so no unstated symmetry is used to assert the intended lower bound.

\textbf{Completion rate estimate: 100\%.}

\end{document}
